\section*{Chapitre \ref{ch:b1:p-block-16-18} --- Le bloc p, II~: oxygène, halogènes et gaz nobles}

\begin{solution}{exo:b1:p-block-16-18:1}
\ce{H2S} $-$II~; \ce{S8} 0~; \ce{SO2} $+$IV~; \ce{SO3^2-} $+$IV~; \ce{SO4^2-}
$+$VI~; \ce{S2O3^2-} $+$II en moyenne~; \ce{H2S2O7} $+$VI.
\end{solution}

\begin{solution}{exo:b1:p-block-16-18:2}
Un \omterm{def:b1:p-block-16-18:halogen}{halogène} oxyde les halogénures situés au-dessous de lui~:
\ce{Cl2 + 2I- -> 2Cl- + I2} et
\ce{Cl2 + 2Br- -> 2Cl- + Br2} ont lieu~; le diiode avec le bromure et le
dibrome avec le chlorure ne réagissent pas.
\end{solution}

\begin{solution}{exo:b1:p-block-16-18:3}
\ce{SO2}~: S porteur d'un \omterm{def:b1:lewis-resonance-vsepr:lewis}{doublet non liant}, une liaison S=O et une
liaison \ce{S+-O-} dans chacune des deux \omterm{def:b1:lewis-resonance-vsepr:resonance}{formules mésomères} (ou deux
S=O avec un octet étendu)~: coudée.
\ce{SO3}~: trigonale plane, trois liaisons S--O équivalentes. \ce{O3}~:
O central porteur d'un \omterm{def:b1:lewis-resonance-vsepr:lewis}{doublet non liant}, une liaison O=O et une liaison
O--O$^-$, deux \omterm{def:b1:lewis-resonance-vsepr:resonance}{formules mésomères}~: coudée.
\end{solution}

\begin{solution}{exo:b1:p-block-16-18:4}
\ce{SO2 + H2O <=> H2SO3}~; \ce{SO2 + 2NaOH -> Na2SO3 + H2O}~;
\ce{SO3 + H2O -> H2SO4}~; \ce{SO3 + 2NaOH -> Na2SO4 + H2O}.
\end{solution}

\begin{solution}{exo:b1:p-block-16-18:5}
HF~: $x^2/(0.10 - x) = 10^{-3.16}$, $x = \qty{8.0e-3}{mol/L}$, pH 2.10.
HCl~:
pH 1.00. La liaison H--F est forte et l'ion fluorure est fortement lié
par \omterm{def:b1:intermolecular-forces-solvents:hbond}{liaisons hydrogène}~: HF est un \omterm{def:b1:predominant-reaction:strong-acid}{acide faible}.
% ledger: dfg:HF_ao, dfg:F-_ao
\end{solution}

\begin{solution}{exo:b1:p-block-16-18:6}
\ce{BrF5}~: cinq liaisons et un \omterm{def:b1:lewis-resonance-vsepr:lewis}{doublet non liant}, pyramide à base
carrée. \ce{IF7}~: bipyramide pentagonale. \ce{ICl4-}~: quatre liaisons
et deux \omterm{def:b1:lewis-resonance-vsepr:lewis}{doublets non liants}, plan carré. \ce{XeO3}~: trois liaisons et
un \omterm{def:b1:lewis-resonance-vsepr:lewis}{doublet non liant}, pyramide trigonale.
\end{solution}

\begin{solution}{exo:b1:p-block-16-18:7}
$\log K = 2(1.36 - 0.53)/0.059 = 28$. Le diiode libéré forme avec
l'amidon le \omterm{def:b1:complexation:complex}{complexe} bleu foncé.
\end{solution}

\begin{solution}{exo:b1:p-block-16-18:8}
$E^\circ(\ce{O3}/\ce{O2}) = 2.07 > 0.53$~V~: \ce{O3 + 2I- + 2H+ -> O2 + I2
+ H2O}. Le diiode formé est titré par le thiosulfate
(\cref{exo:b1:nernst:7})~: un \ce{I2} par \ce{O3}.
\end{solution}

\begin{solution}{exo:b1:p-block-16-18:9}
\ce{C12H22O11 -> 12C + 11H2O}~; $M = \qty{342.0}{g/mol}$, $10.0/342.0 =
\qty{0.0292}{mol}$, carbone~: $0.0292 \times 12 \times 12.0 = \qty{4.21}{g}$.
\end{solution}

\begin{solution}{exo:b1:p-block-16-18:10}
$x(0.10 + x)/(0.10 - x) = 10^{-1.99}$~: $x = \qty{8.6e-3}{mol/L}$,
$[\ce{H3O+}] = 0.1086$, $\mathrm{pH} = 0.96$ au lieu de 1.00~: la seconde
acidité augmente $[\ce{H3O+}]$ d'environ \qty{9}{\percent}.
\end{solution}

\begin{solution}{exo:b1:p-block-16-18:11}
\ce{F2}/\ce{F-} (2.89~V) est situé au-dessus de tous les autres
couples~: aucun \omterm{def:b1:nernst:couple}{oxydant} ne peut arracher l'électron de l'ion fluorure.
Lors d'une électrolyse en solution aqueuse, c'est l'eau, oxydée à un
potentiel bien plus bas (\ce{O2}/\ce{H2O}, 1.23~V), qui réagirait~; un
fluorure fondu ne contient pas d'eau.
\end{solution}

\begin{solution}{exo:b1:p-block-16-18:12}
$\Delta_r G^\circ = -51.4 - (16.40 - 51.57) = \qty{-16.2}{kJ/mol}$,
$\log K = 16.2/5.708 = 2.84$, $K = 7 \times 10^2$.
% ledger: dfg:I2_ao, dfg:I-_ao, dfg:I3-_ao
\end{solution}

\begin{solution}{pb:b1:p-block-16-18:1}
\textbf{1.} 0, $+$IV, $+$VI, $+$VI.
\textbf{2.} \ce{S + O2 -> SO2}.
\textbf{3.} S porteur de deux domaines liants et d'un \omterm{def:b1:lewis-resonance-vsepr:lewis}{doublet non
liant}~: molécule coudée, environ
\ang{120}.
\textbf{4.} $10^6/32.1 = \qty{3.12e4}{mol}$.
\textbf{5.} \qty{3.115e4}{mol} de \ce{O2}~: $V = 3.115 \times 10^4 \times
8.314 \times 298.15/10^5 = \qty{772}{m^3}$, air~: $772/0.21 = \qty{3.7e3}{m^3}$.
\textbf{6.} La combustion est très exothermique~; le \omterm{def:b1:elementary-steps:catalyst}{catalyseur}
fonctionne dans un domaine de température limité, et l'oxydation de
\ce{SO2}, exothermique elle aussi, est moins complète à chaud.
\textbf{7.} \ce{2SO2 + O2 <=> 2SO3}.
\textbf{8.} Une grande constante d'équilibre ne dit rien de la vitesse~:
à température ambiante, la réaction n'avance pas.
\textbf{9.} Il offre un chemin d'\omterm{def:b1:rate-laws:activation-energy}{énergie d'activation} plus faible et est
régénéré~: c'est un \omterm{def:b1:elementary-steps:catalyst}{catalyseur} hétérogène.
\textbf{10.} Trigonale plane.
\textbf{11.} Davantage de dioxygène maintient le quotient
$Q = p_{\ce{SO3}}^2/(p_{\ce{SO2}}^2 p_{\ce{O2}})$ plus bas~: la
conversion
de \ce{SO2} va plus loin.
\textbf{12.} \qty{0.3}{\percent}.
\textbf{13.} \ce{SO3 + H2SO4 -> H2S2O7}.
\textbf{14.} Il forme un brouillard de fines gouttelettes d'acide qui
traverse l'absorbeur.
\textbf{15.} \ce{H2S2O7 + H2O -> 2H2SO4}.
\textbf{16.} \ce{2S + 3O2 + 2H2O -> 2H2SO4}.
\textbf{17.} Une molécule d'eau par atome de soufre~: $\num{3.115e4} \times 18.0 =
\qty{561}{kg}$.
\textbf{18.} La dilution dégage beaucoup de chaleur~; de l'eau versée
sur l'acide bout aussitôt et projette de l'acide.
\textbf{19.} $x(0.050 + x)/(0.050 - x) = 10^{-1.99}$~: $x = \qty{7.5e-3}{mol/L}$,
$[\ce{H3O+}] = \qty{0.0575}{mol/L}$, $\mathrm{pH} = 1.24$.
\textbf{20.} $3.12 \times 10^4 \times 0.995 \times 98.1 = \qty{3.04}{t}$.
\textbf{21.} $84 \times 98.1/32.1 = \qty{257}{Mt}$.
\textbf{22.} $98.1/32.1 =$ \textbf{\qty{3.06}{t} d'acide sulfurique par
tonne de soufre}.
\end{solution}
